\documentclass[10pt]{article}
% TMLR style (https://github.com/JmlrOrg/tmlr-style-file).
% [preprint] shows the author names and drops the "Under review as submission to TMLR" header.
\usepackage[preprint]{tmlr}

\usepackage{hyperref}
\usepackage{url}
\usepackage{graphicx}
\usepackage{booktabs}

\title{EvoSym: Evolvable Neuro-Symbolic Intelligence}

% Separate authors with \AND.
\author{\name Gian Seifert \\
      \addr University of St.\ Gallen
      \AND
      \name Cyril Gabriele \\
      \addr University of St.\ Gallen
      \AND
      \name Filipp Rechsteiner \\
      \addr University of St.\ Gallen}

\begin{document}

\maketitle

\begin{abstract}
A question such as ``Where can a train take me from Bern without stopping?'' needs facts from several messy SBB Open Data exports, rules that combine them, and an explainable answer. With EvoSym we built a FrameX knowledge base that turns nine of these exports into 443,975 facts under an open-world assumption. Eight rules derive new facts, and every derived answer traces back to its rule and source facts. All 24 tests pass, 14 of the 18 test queries return exactly the reference answer, and the other four differ for known data reasons. The median query takes 0.22\,ms.

Cyril Gabriele served as the developer, Filipp Rechsteiner as the domain expert, and Gian Seifert as the auditor.
\end{abstract}

\section{Introduction}
\label{sec:introduction}

% Task: Hackathon 1 task sheet and Hybrid AI Studio slides 03, "Hackathon 1: Derive From Knowledge".
% Numbers: Sections ingestion, ontology, results (performance) of this report.
Large language models answer questions about train travel fluently, but they have no explicit notion of facts, may hallucinate and cannot show why an answer holds. A symbolic system stores facts explicitly, derives new facts with rules and can explain every conclusion. Hackathon~1 of the Hybrid AI course asked for such a system: a FrameX~\citep{framex} program that ingests messy open data of the Swiss Federal Railways (SBB)~\citep{sbbdata}, models it with an ontology and rules, derives explainable new facts and answers multi-hop questions about Swiss passenger rail travel without an LLM or internet access after initialisation.

This report presents EvoSym, the team's solution. EvoSym turns nine SBB Open Data exports into 443,975 FrameX facts under an open-world assumption, derives new facts with eight rules and answers the 18 test queries provided by the teaching assistants. Every derived answer traces back to its rule and source facts. On the submitted data, 14 of the 18 queries return exactly the reference answer, the other four differ for known data reasons, and the median query takes 0.22\,ms. Section~\ref{sec:methods} describes the data ingestion, the ontology and rules, the Python integration, the world assumption and the evaluation. Section~\ref{sec:results} reports correctness and performance, and Section~\ref{sec:discussion} discusses the limits.

\section{Methods}
\label{sec:methods}

\subsection{Data ingestion}
\label{sec:ingestion}


As the domain expert, the mandate was to understand the SBB data and its limits and to turn the questions into requirements for the facts. This covered choosing the datasets, defining the scope and writing the first version of the ingestion script, documented in \texttt{docs/hackathon01/DATA\_INGESTION.md}. The developer later moved the script into \texttt{src/adapters/sbb/} and added a Pydantic validation.

The datasets were chosen according to the TA queries. Of the 60 datasets on SBB Open Data~\citep{sbbdata}, the nine that the queries draw on were used; most others describe infrastructure, real estate or ticket products. They were downloaded through the Explore API and stored unchanged, so every later step runs offline. A Didok service point was kept only if it lies in Switzerland, is a stop point and is served by train or rack railway: 1,773 of 60,162 points. The filter uses Didok's country, not the number prefix, because Jestetten and Lottstetten carry Swiss numbers but lie in Germany. The datasets store the stop-point number as an integer, as text or as a float (\texttt{8502113.0}). Each value was converted to one identifier, \texttt{sp\_<number>}, so a query can join facts across datasets.

The train runs needed the most cleaning. Of 69,854 stop records, 1,350 cancelled stops, 32 pass-throughs and 3,930 stops abroad were dropped. The stops of each run were sorted by time, and consecutive stops were linked with \texttt{nextStop}, on which the rule for direct connections builds (Section~\ref{sec:ontology}).

The facts contain only what the data states; derived properties such as long platforms are left to the rules. The nine files hold 443,975 facts, and each file header records the source, the download hash, the counts and the dropped rows.

\paragraph{Data limitations.}
Six WiFi rows have no stop-point number (Section~\ref{sec:world}). The train export holds only the previous day, so three queries differ from the reference. Three planned halls at Langenthal have no building name, but they were kept because stop point, canton and status are present, which explains the fourth difference (Table~\ref{tab:reference}).

\subsection{Ontology and derivation rules}
\label{sec:ontology}

% Mandate: docs/exercise/02_week/02_HybridAI_Studio.pdf, PDF pages 8 and 9.
% Contributions: Cyril's commits 27cb52d and a2763cc; docs/hackathon01/roles/hackathon01/cyril.md.
My mandate as the developer was to own the software architecture, integrate FrameX, and keep the code runnable, structured and extensible. I separated the ontology, source facts, derivation rules and queries into distinct files. I defined a class hierarchy for places, facilities, train runs and stop events, with typed properties connecting them. I used shared station identifiers to connect facts across the SBB datasets.

Separating these responsibilities makes the system easier to inspect and maintain. Source facts record the input evidence, while FrameX rules express deductions and provide proofs for derived answers. If an auditor finds an incorrect answer, this separation helps isolate whether the cause lies in the source data, its translation into facts, a derivation rule or the query.

This implementation covers only part of the more detailed ontology we designed during FlightMode. Following Götz's advice to anticipate future hackathons, we included finer distinctions beyond the needs of the current queries. In theory, these distinctions could give an LLM agent a richer semantic understanding of the domain.

For this hackathon, I implemented eight FrameX rules for long platforms, waiting hall availability, junctions, busy stations by year, train categories serving a station, stations served by long distance trains, direct connections between consecutive stops, and train and tram interchanges. For example, I required two distinct infrastructure lines for a junction and a length above 320\,m for a long platform. For direct connections, I linked consecutive stop events within the same train run to avoid inferring a direct connection across intermediate stops.

\subsection{Python integration}
\label{sec:integration}

I built the Python client in \texttt{src/framex.py} using only the standard library to communicate with the local FrameX process through JSON messages. I added a runner that combines the ontology, rules and either sample or ingested facts, loading large programs in chunks. Through \texttt{src/main.py}, I exposed all 18 queries, explanations and schema validation. I also wrote a small sample of facts so the team could exercise every query without downloading SBB data, matching my completion criterion.

For future hackathons, a developer can add rules and queries while reusing the existing FrameX client and runner, extending the ontology where necessary. Adding datasets requires changes to the ingestion pipeline and the runner's dataset list. Splitting the SBB adapter into modules for downloading, dataset conversion and output handling would make such extensions more manageable.

\paragraph{Implementation difficulties.}
The tight time limit left little room to refine the architecture. In particular, the SBB ingestion adapter takes on too many responsibilities, and I recognised this too late to refactor it.

\subsection{World assumption}
\label{sec:world}

Every program starts with \texttt{world open.}~\citep{framex}. A missing fact therefore means \emph{unknown}, not \emph{false}. We chose the open world because the SBB exports~\citep{sbbdata} are real-world data and none of them claims to be complete. For example, six rows of the WiFi export have no stop-point number, among them St.~Gallen. These stations get no \texttt{hasWifi} fact, so a closed world would claim that St.~Gallen has no WiFi. The adapter writes positive facts only, and an empty source value writes no fact. Query~1.7 (long-distance stations without WiFi) therefore returns \texttt{unknown}, as the reference expects.

\subsection{Evaluation}
\label{sec:evaluation}

The evaluation checks two things: that every answer follows from the rules and the source data, and that the system answers \texttt{unknown} when the evidence is missing. We started with the six questions from our FlightMode concept and switched to the 18 TA queries once they were released. The queries run word for word, so the team adopted the TA vocabulary (\texttt{sp\_<UIC>}, \texttt{atStopPoint}, \ldots) as its naming contract. All tests use Python's \texttt{unittest} and need no extra dependency. Table~\ref{tab:tests} lists the four kinds of tests.

\begin{table}[h!]
\centering
\caption{Test suite. Sample tests run on a small hand-written fact file and need no download. Real-data tests run with \texttt{RUN\_SBB\_TESTS=1}.}
\label{tab:tests}
\small
\begin{tabular}{@{}p{0.2\linewidth}p{0.47\linewidth}p{0.25\linewidth}@{}}
\toprule
Test & What it checks & Data \\
\midrule
Query answers & All 18 queries return the expected rows; one subtest per query & Sample and real data \\
Near misses & Cases just outside a rule must stay \texttt{unknown} & Sample \\
Explanations & A derived fact traces back to a rule and an asserted source fact & Sample (8 rules), real data (3 facts) \\
Counterfactuals & Changing one supporting fact removes the conclusion & Sample (8 rules, 10 cases) \\
\bottomrule
\end{tabular}
\end{table}

\paragraph{Independent expected answers.}
\texttt{scripts/build\_expected.py} computes the expected answer to each query directly from the raw JSON exports in plain Python. It uses neither FrameX nor the ingestion adapter nor the rules, so a bug in those cannot also hide in the expected answer. Before any rule existed, we fixed how to read each question by looking at the raw data:
\begin{itemize}\setlength{\itemsep}{0pt}\setlength{\parskip}{0pt}
  \item Scope: Swiss service points with train or rack-railway service. Without this filter an emergency stop inside the Gotthard Base Tunnel counts as a station in Graub\"unden.
  \item Daily passengers are the average over all days (DTV), and thresholds are strict. With the weekday average, Biel/Bienne would also pass query~1.5.
  \item A train stopped at a station only if it neither passed through nor was cancelled.
  \item A waiting hall counts if it exists or is planned for demolition, but not if it is only planned. We first counted existing halls only and changed this after the reference counted Landquart, whose hall still stands.
  \item Long-distance categories are EC, IC, ICE, IR, NJ, RJ, RJX and TGV.
\end{itemize}
The train-run export holds only the previous operating day and is replaced daily. The expected answers therefore store SHA-256 hashes of the downloads they were built from, and the real-data tests refuse to compare against any other download.

\paragraph{Near misses and explanations.}
The sample contains cases that must stay \texttt{unknown}: a platform of exactly 320\,m, a station on a single line, a hall that is only planned, a passenger count just below the threshold, a station with no WiFi row, and a train that stops in Z\"urich between Bern and Chur, so Bern is not non-stop to Chur. For one derived fact per rule, \texttt{explain} must return a proof that contains both a rule step and an asserted source fact.

\paragraph{Counterfactuals.}
In FlightMode we asked whether the system could derive a conclusion that does not depend on its evidence, for example that Chur has WiFi when it does not. For each of the eight rules, a test takes a conclusion that is true on the sample, changes one supporting fact, reloads the program and requires the conclusion to become \texttt{unknown}. Examples are removing Chur's WiFi fact, shortening a 321\,m platform to 320\,m and changing an IC run to an S-Bahn. Two guards keep the test honest. First, the conclusion must be \texttt{true} before the edit: we found that query~1.7 also returns \texttt{unknown} on a program with no facts at all, so a broken rule would otherwise pass. Second, each edit must match exactly one place in the source, so a typo cannot silently change nothing. This guard caught one of our own edits, whose text also appeared in \texttt{rules.fx}. On the real data, removing Chur's WiFi fact turned ``Chur has WiFi'' from \texttt{true} to \texttt{unknown} and left Chur's junction status unchanged.

\paragraph{Testing the tests.}
We checked that the tests can fail. On an early synthetic knowledge base, deleting the Langenthal halls and the TGV stops failed exactly queries~2.4 and~3.4. Removing the condition \texttt{?A != ?B} from the junction rule made a station on one line a junction, and the counterfactual test failed.

\section{Results}
\label{sec:results}

\paragraph{Correctness.}
All 24 tests pass. On 2 October 2026 we rebuilt the expected answers from a download for operating day 30 September 2026, and the real-data tests passed in 7.1\,s. On the submitted data (operating day 28 September 2026), 14 of the 18 queries return exactly the TA reference. The other four differ for known reasons (Table~\ref{tab:reference}).

\begin{table}[t]
\centering
\caption{Queries whose answers differ from the TA reference.}
\label{tab:reference}
\small
\begin{tabular}{@{}lp{0.8\linewidth}@{}}
\toprule
Query & Reason \\
\midrule
1.4, 1.6, 3.3 & Train data: our snapshot covers 28 September 2026, the reference 27 September 2026. The export only holds the previous day, so we cannot download the reference day. \\
2.4 & We return three more planned halls at Langenthal. Their source rows state canton BE and status \texttt{PROJEKTIERT NEU}, so we keep them. \\
\bottomrule
\end{tabular}
\end{table}

\paragraph{Performance.}
Measured on an Apple Silicon laptop, the engine starts in 3.9\,ms. Loading the 443,975 source facts and evaluating the rules takes 5.9\,s and yields 537,571 facts, 93,539 of them derived (226,112 rule firings in three rounds). Because the engine derives everything at load time, a query only looks up facts: the median query takes 0.22\,ms. The weak point is memory, with 1,353\,MB for the engine process.

\section{Discussion}
\label{sec:discussion}

% TODO (team): main discussion.

The evaluation has four limits.
\begin{itemize}\setlength{\itemsep}{0pt}\setlength{\parskip}{0pt}
  \item The expected answers repeat our reading of each question in Python. A misread question is wrong in both places and still passes. Only the comparison with the TA reference guards against this.
  \item The tests catch wrong reasoning, not wrong data. If SBB listed a station with WiFi that has none, the system would correctly derive \texttt{true} from a false source.
  \item The explanation tests check that a proof contains a rule and an asserted fact, not that it uses the right rule and the right facts.
  \item Answers about trains describe a single operating day. After each download the expected answers must be rebuilt, and the reference answers of another day can only guide the meaning of a query, not confirm our answers.
\end{itemize}

\bibliography{references}
\bibliographystyle{tmlr}

\appendix
\section{Appendix}

\subsection{Use of coding agents}
\label{app:agents}

Coding agents were used only after FlightMode. Gian used Claude Code with the models Claude Opus 5.5 and Claude Fable 5.1 for the tests, the expected-answer script, the performance report, the README and the architecture decision records. Every result was checked by running it, and the expected answers were written independently of the code they test.

Cyril used Codex with GPT-6 Sol at medium reasoning effort for coding, and briefly used Astra 6 at medium reasoning effort for the SBB ingestor.

Filipp used Claude Code with the model Opus 5.5 for coding.

\end{document}
